\documentclass[11pt]{article}

\usepackage[margin=1in]{geometry}
\usepackage{graphicx}
\usepackage{booktabs}
\usepackage{amsmath,amssymb}
\usepackage{xcolor}
\usepackage{hyperref}
\usepackage{natbib}
\bibliographystyle{unsrtnat}

\title{Deep-learning-guided design of 5$^\prime$UTRs for efficient
mRNA-delivered gene editing}

\author{Anonymous Authors\\Target venue: Nature Biotechnology}
\date{}

\begin{document}
\maketitle

\begin{abstract}
Messenger RNA therapeutics depend critically on efficient translation of the
encoded protein, yet the 5$^\prime$ untranslated region (5$^\prime$UTR) has
remained one of the least systematically optimized components of the mRNA
sequence. Most deep-learning models of 5$^\prime$UTR activity are trained on
fixed-length reporter libraries in a single cell type, and their ability to
guide \emph{functional} design of therapeutically relevant cargos such as
mRNA-delivered gene editors is largely untested. Here we quantify 5$^\prime$UTR
activity at scale across three human cell types by polysome-profiling massively
parallel reporter assays (MPRAs), train convolutional and k-mer models of mean
ribosome load (MRL), and apply gradient-based Fast SeqProp and generative Deep
Exploration Networks (DENs) to design both 50- and 25-nucleotide 5$^\prime$UTRs
\emph{de novo}. MRL is tightly correlated across cell lines
($r^{2} = 0.837\text{--}0.870$ between HEK293T and T cells, $0.847\text{--}0.871$
between HEK293T and HepG2) and nearly as reproducible as replicates from the
same cell line. The HEK293T-trained Optimus 5-Prime model reaches
$r^{2} = 0.937$ on held-out HEK293T sequences and $r^{2} = 0.841$ and $0.840$
on T cell and HepG2 data. Existing Optimus models transfer poorly to a new
fully-variable 25~nt library ($r^{2} = 0.564$ and $0.600$), motivating a
retrained Optimus 5-Prime(25) that reaches $r^{2} = 0.806$ on held-out data.
When coupled to mRNA-encoded megaTAL gene editors in K562 cells, DEN-designed
5$^\prime$UTRs drive absolute editing efficiencies exceeding 40\% for
\emph{TGFBR2} and 80\% for \emph{PDCD1} at 2~pmol mRNA; the best DEN design
reaches 55.6\% on \emph{TGFBR2}, 18--33\% higher than the previous best natural
control LAMA5. Cargo- and locus-specific deviations --- the same DEN design
yields only 80.8\% on \emph{PDCD1} while VAT1 reaches 91.4\% --- indicate that
predicted MRL is a useful but imperfect proxy for editing efficiency,
reinforcing the need for empirical validation in the deployment context.
Together, these results show that model-guided 5$^\prime$UTR design is a
practical tuning knob for mRNA therapeutics, and supply a benchmark for
evaluating future generative approaches.
\end{abstract}

\section{Introduction}

Messenger RNA has moved rapidly from a niche laboratory reagent to a central
modality for vaccines, protein replacement, and \emph{in situ} genetic
medicine. Nucleoside-modified, lipid-nanoparticle-formulated mRNAs are the
basis of the first approved SARS-CoV-2 vaccines~\citep{thomas2021bnt162b2,
karik{\'o}2013pseudouridine}, underpin individualized cancer vaccines now
entering durable-response regimes~\citep{sahin2026tnbcvaccine}, and are being
formulated for targeted delivery to lymph nodes and other
tissues~\citep{kim2026lnp_utr,wisitrasameewong2022mrnaregen}. Because every
such therapy depends on the amount of protein produced per delivered mRNA
molecule, the primary sequence of the mRNA itself --- and in particular its
5$^\prime$ untranslated region --- is directly tied to therapeutic potency.
Yet the 5$^\prime$UTR has remained surprisingly under-engineered relative to
lipid chemistry and codon usage.

The 5$^\prime$UTR determines how the ribosome engages the mRNA and how it
selects the start codon. Classical determinants include the Kozak context at
the initiation AUG, the presence and strength of upstream open reading frames
(uORFs), local secondary structure, and G/C composition along the scanning
tract~\citep{chaldebas2026utrvariants,stettler2026kozak,
balasov2025kozaknoncanonical,shikalov2021ddx3}. These features interact
non-linearly, and even modest changes to the 5$^\prime$UTR can alter protein
output by orders of magnitude. The past several years have seen massively
parallel reporter assays (MPRAs) coupled to polysome profiling emerge as a
practical way to quantify 5$^\prime$UTR activity at the library
scale~\citep{reimaopinto2025danio5p,strayer2024naptrap}, yielding training
data for supervised neural networks that predict mean ribosome load (MRL)
directly from the sequence. However, reporter libraries have typically been
profiled in a single cell line, usually HEK293T, and it is not obvious that
models trained in this way generalize to the cells most relevant for mRNA
therapeutics, such as primary T lymphocytes or liver-derived cells.

In parallel, generative modelling of biological sequences has matured.
Gradient-based activation maximization over one-hot relaxations, as in
Fast SeqProp~\citep{linder2021fastseqprop}, produces high-scoring
sequences but suffers from local minima and a tendency to collapse
onto low-diversity motifs. Deep Exploration Networks
(DENs)~\citep{linder2020den} overcome both issues by training a
generator to maximize a fitness predictor while simultaneously
penalizing pairwise similarity between outputs, and they can be
combined with variational autoencoders to keep generated sequences
inside the data manifold. Similar ideas have powered recent
generative systems for promoters~\citep{zhao2024promotergan,
lei2025promoterdiffusion} and m$^{1}\Psi$-modified 5$^\prime$UTRs for
vaccines~\citep{tang2023smart5utr}. Whether these design strategies
translate into \emph{functional} gains when the payload is a
clinically relevant cargo --- rather than a fluorescent reporter ---
is largely an open question.

Gene editing by mRNA-delivered nucleases is a particularly demanding
test bed. Editing efficiency depends non-linearly on transient enzyme
concentration, and the nucleases themselves are large, structured
proteins that can be sensitive to translation context. megaTAL
nucleases~\citep{mcmurrough2018laglidadg,laforet2019endonuclease,
gwiazda2016tcellediting}, which fuse engineered LAGLIDADG homing
endonuclease scaffolds to TALE DNA-binding arrays, have been used to
disrupt clinically relevant loci such as \emph{TGFBR2} in
tumor-infiltrating lymphocytes~\citep{fix2022tgfbr2} and \emph{PDCD1}
in T-cell
products~\citep{mohammad2025pd1}. Electroporated synthetic mRNA is
a standard delivery mode for these
edits~\citep{schwarze2021electroporation,gwiazda2016tcellediting},
which makes the translational tuning of the 5$^\prime$UTR directly
consequential for editing outcomes.

In this work, we connect MPRA-based modelling to mRNA-delivered
gene editing in a closed loop. We profile randomized 5$^\prime$UTR
libraries in three human cell types by polysome-profiling MPRA, train
convolutional and k-mer models on the resulting data, and use them to
guide the design of both 50- and 25-nucleotide 5$^\prime$UTRs via Fast
SeqProp and DENs. We then place designed 5$^\prime$UTRs upstream of
megaTAL nucleases and read out editing at two independent loci in two
cell lines. Three findings structure the paper: first, MRL is highly
correlated across cell types and therefore transferable; second, a
model matched to the library length (Optimus 5-Prime(25)) is needed
for fully-variable 25~nt 5$^\prime$UTRs; and third, generatively
designed 5$^\prime$UTRs support therapeutically meaningful editing
(above 40\% for \emph{TGFBR2} and above 80\% for \emph{PDCD1} at
2~pmol mRNA), but with cargo- and locus-specific deviations that argue
for empirical validation alongside any purely \emph{in silico}
pipeline.

\section{Related Work}

\paragraph{Deep-learning models of 5$^\prime$UTR translation.}
MPRA-coupled deep models have uncovered the sequence grammar that
controls ribosome recruitment, including 5$^\prime$UTR length,
start-codon context, upstream AUGs, and short
motifs~\citep{reimaopinto2025danio5p,strayer2024naptrap,
chaldebas2026utrvariants}. These studies typically profile a single
cell type, so the question of whether the learned grammar is
portable across cell lines has remained open. Our cross-cell-line
MPRA design directly addresses this point.

\paragraph{Generative sequence design.}
Gradient ascent over continuous one-hot relaxations
(Fast SeqProp) and generative networks with explicit diversity
penalties (DENs) are the current workhorses of cis-regulatory
design~\citep{linder2020den,linder2021fastseqprop}. Complementary
GAN- and diffusion-based generators have been applied to
synthetic promoters~\citep{zhao2024promotergan,
lei2025promoterdiffusion}, and Smart5UTR targets
m$^{1}\Psi$-modified 5$^\prime$UTRs for mRNA
vaccines~\citep{tang2023smart5utr}. Our contribution here is not a
new generator architecture, but a functional evaluation that
couples existing Fast SeqProp and DEN machinery to an mRNA-delivered
editing readout.

\paragraph{mRNA therapeutics.}
Nucleoside-modified, LNP-formulated mRNA underpins approved vaccines
and clinical-stage individualized cancer
therapies~\citep{thomas2021bnt162b2,sahin2026tnbcvaccine,
karik{\'o}2013pseudouridine}. Recent efforts couple ionizable lipid
engineering with UTR tuning to achieve potency at low
doses~\citep{kim2026lnp_utr} and to stabilize mRNAs against host
translational suppression~\citep{phua2026ns1translation}. Our 5$^\prime$UTR
designs are complementary to, and composable with, such formulation
advances.

\paragraph{megaTAL-based gene editing.}
megaTALs combine engineered LAGLIDADG homing endonucleases with TALE
DNA-binding arrays and edit primary human T cells at high
efficiency~\citep{gwiazda2016tcellediting,mcmurrough2018laglidadg,
laforet2019endonuclease}. \emph{TGFBR2} knockout liberates
tumor-infiltrating lymphocytes from TGF-$\beta$
suppression~\citep{fix2022tgfbr2}, and \emph{PDCD1} disruption is
being pursued to counter T-cell exhaustion in CAR-T and TIL
products~\citep{mohammad2025pd1}. Because mRNA electroporation is a
standard delivery route for these edits, every
percentage-point improvement in translation can matter clinically.

\section{Methods}

\subsection{Polysome-profiling MPRA in three cell types}

We assembled two synthetic 5$^\prime$UTR libraries upstream of a common
reporter: a fixed-length random 50-nucleotide library
(primer Bri036\_T7\_N50\_ATG) and a random-end 25-nucleotide library
(primer Bri043\_T7\_N25\_ATG). Libraries were assembled with 200~ng of
backbone and 10~pmol of primer using NEBuilder HiFi DNA Assembly
Master Mix in 20~$\mu$L reactions incubated at 50$^\circ$C for 1~hour;
products were purified with DNA Clean \& Concentrator-5 columns (Zymo)
with elution in 7~$\mu$L and transformed into NEB 5-alpha cells in two
parallel reactions of 3.5~$\mu$L DNA in 35~$\mu$L cells.

Libraries were delivered as mRNA and polysome-profiled in three human cell
types: HEK293T, primary T cells, and HepG2. HepG2 cells were cultured in
EMEM supplemented with 10\% FBS. Primary T cells were plated in T-cell
growth medium (TCGM) at $10^{6}$~cells~mL$^{-1}$ and incubated at
$37^\circ$C and 5\%~CO$_2$; 8~hours after transfection, cycloheximide was
added to a final concentration of 100~$\mu$g~mL$^{-1}$ and cells were
incubated for an additional 5~minutes before harvesting. Cells were
pelleted at 1500~rpm for 5~min, resuspended in 300~$\mu$L of cold lysis
buffer, and incubated on ice for 10~minutes. Lysates were spun at
16{,}000~rpm for 5~min at 4$^\circ$C, treated with 1.5~$\mu$L of
1~U~$\mu$L$^{-1}$ DNase, incubated on ice for 30~minutes, and stored at
$-80^\circ$C.

Extracted RNA was eluted in 11~$\mu$L of RNase-free water. Reverse
transcription was initiated by mixing 10.5~$\mu$L of purified RNA with
2~$\mu$L of 2~$\mu$M RT primer (carrying a 10-nucleotide UMI) and
2~$\mu$L of 10~mM dNTPs (NEB), followed by 5~minutes at 65$^\circ$C and
at least 1~minute on ice. We then added 4~$\mu$L Maxima RT buffer,
0.5~$\mu$L Superase-In, and 1~$\mu$L Maxima RT RNaseH minus enzyme
(Thermo Fisher), incubated at 50$^\circ$C for 15~minutes and
85$^\circ$C for 10~minutes, and returned to ice. Residual RNA was
degraded with 1~$\mu$L RNase (Roche, bovine pancreas, DNase free) plus
1~$\mu$L RNase H (NEB) at 37$^\circ$C for 15~minutes, after which the
product was cleaned with 3$\times$ KAPA beads (Roche) and resuspended
in 25~$\mu$L of water. Template switching used 8~$\mu$L Maxima RT
buffer, 6~$\mu$L 50\% PEG-8000, 0.5~$\mu$L Superase-In, 1~$\mu$L
Maxima RT RNaseH minus enzyme, and 0.5~$\mu$L of 10~$\mu$M template
switching oligo (AAGCAGTGGTATCAACGCAGAGTACATrGrGrG); another RNase/
RNase H digest and a 3$\times$ KAPA bead cleanup into 20~$\mu$L of
water followed. Amplification used KAPA HiFi master mix with forward
and reverse primers carrying 8-nt index barcodes specific to each
polysome fraction, stopped before saturation, and purified by gel
extraction. Final libraries were pooled, cleaned with Ampure XP beads,
normalized to 16~pM, and sequenced on an Illumina MiSeq with
25\% PhiX spike-in.

Read counts were tallied per 5$^\prime$UTR variant and summarized as
mean ribosome load (MRL). Analyses focused on high-coverage subsets:
the top 20{,}000 sequences by read coverage for the 50~nt library and
the top 2{,}000 for the 25~nt library, where inter-replicate
correlations were highest.

\subsection{MRL models: Optimus 5-Prime and Optimus 5-Prime(25)}

We trained a convolutional neural network (Optimus 5-Prime) on the
fixed-end 50~nt HEK293T MPRA data to predict per-sequence MRL, and
retrained the architecture on the random-end 25~nt data to obtain
Optimus 5-Prime(25). For validation of gradient-based designs, we
trained companion linear k-mer predictors: the 50~nt k-mer model
reached Pearson $r = 0.5213$ on its test set, and the 25~nt k-mer
model reached Pearson $r = 0.4094$ on a held-out set of 11{,}349
sequences with no uAUG and $>$200 reads. All neural network training
was performed in Python~3 with Keras~2.2 and TensorFlow~1.15.

\subsection{Sequence design: Fast SeqProp and Deep Exploration Networks}

Fast SeqProp~\citep{linder2021fastseqprop} designs sequences by
gradient-based activation maximization on a relaxed one-hot
representation; we ran SeqProp both with and without VAE regularization.
VAE training sets were constructed from MPRA data by first filtering
on coverage ($>$2{,}000 reads for the 50~nt library, $>$500 reads for
the 25~nt library) and then drawing 5{,}000 training and 1{,}000 test
sequences from the top 10{,}000 by measured MRL
(approximately the top 25\%).

DENs~\citep{linder2020den} map a 100-dimensional Gaussian latent vector
through a generator network $G_\theta$ to a $50\times4$ or $25\times4$
continuous logit that, after softmax or sampling, yields a candidate
5$^\prime$UTR (Figure~\ref{fig:den_schematic}). The training
objective combined a fitness term (the negative Optimus 5-Prime-predicted
MRL) with a pairwise similarity penalty across the generated batch to
explicitly encourage diversity; an optional VAE term maintained the
likelihood ratio of generated sequences under the data
distribution~\citep{linder2020den}. In an ``inverse regression'' mode
the generator received an additional input specifying a target MRL,
and we drew four 50~nt sequences with predicted MRLs of 2, 3.5, 5,
and 6.5 to calibrate the model's dynamic range.

\subsection{mRNA-delivered megaTAL editing assay}

Designed 5$^\prime$UTRs were cloned upstream of megaTAL coding sequences
targeting either \emph{TGFBR2} or \emph{PDCD1}. Following DNase treatment
to remove residual template, mRNA was purified with RNase-free Ampure
beads, size- and purity-checked on a Fragment Analyzer (Agilent),
normalized to 500~nM, and stored at $-80^\circ$C. Doses of 0.25--2~pmol
mRNA were electroporated in duplicate into $10^{5}$ K562 or HepG2 cells
per well using a Lonza 4D-Nucleofector 96-well Shuttle. Editing was
read out by amplicon deep sequencing: 1.5~$\mu$L of genomic DNA was
amplified for 30 cycles with gene-specific primers carrying Illumina
overhangs (PCR1), followed by a barcoding step and a MiSeq run.
Editing efficiency was computed as the fraction of reads carrying
indels at the expected cut site.

\subsection{Controls}

For the 50~nt design assays, controls included eight sequences from a
prior polysome-profiling MPRA --- four with low or medium measured MRL
and four drawn from the top 0.02\% by MRL --- plus a panel of natural
5$^\prime$UTRs including LAMA5 and VAT1, previously reported to have
high translation efficiencies comparable to the commonly-used
$\beta$-globin UTR (top $\sim$20\% among $\sim$17{,}000 short endogenous
5$^\prime$UTRs). For 25~nt designs, we additionally selected four
5$^\prime$UTRs from the random-end 25~nt library with MRLs in the top
0.25\%, after filtering out sequences with fewer than 1{,}000 reads or
those containing uATGs.

\section{Results}

\subsection{5$^\prime$UTR activity is correlated across cell types}

We first asked whether 5$^\prime$UTRs rank-order similarly when delivered to
cells with very different proteomes, since any cross-cell-line transfer of a
5$^\prime$UTR model requires that assumption. We sequenced the random 50~nt
library in HEK293T, primary T cells, and HepG2, and compared the resulting
MRLs on the 20{,}000 sequences with the highest read coverage
(Figure~\ref{fig:cross_cell}).

Between-cell-line correlations were nearly as high as replicate correlations
within a single cell line. We observed coefficients of determination of
$r^{2} = 0.837$--$0.870$ between HEK293T and T cells and
$r^{2} = 0.847$--$0.871$ between HEK293T and HepG2, compared with replicate
$r^{2}$ of $0.938$ for HEK293T and $0.814$ for T cells. The fact that
cross-cell-line comparisons approach the ceiling set by biological replicates
within a cell line indicates that the bulk of the variance in 5$^\prime$UTR
activity is cell-type-independent, at least across these three human cell
backgrounds.

Building on this observation, we trained Optimus 5-Prime on the HEK293T MRL
data and asked how well it predicts MRL in the other cell types. On 20{,}000
sequences with the highest coverage held out from training, the HEK293T model
achieved $r^{2} = 0.937$ against HEK293T measurements; against T cell and
HepG2 measurements it still achieved $r^{2} = 0.841$ and $r^{2} = 0.840$,
respectively (Figure~\ref{fig:cross_cell} and Table~\ref{tab:cross_cell}).
This near-linear drop from $r^{2} \approx 0.94$ to $r^{2} \approx 0.84$
closely tracks the inter-cell-line MRL correlation itself, suggesting that
the loss in predictive accuracy across cell types is explained almost
entirely by biological cell-type variability rather than by a failure of the
model. Practically, it means a single HEK293T-trained model is a reasonable
starting point for 5$^\prime$UTR design in T cells or hepatocytes, provided
one accepts the extra noise floor introduced by cell-type differences.

\begin{table}[htbp]
\centering
\caption{Coefficients of determination ($r^{2}$) for mean ribosome load
(MRL) across cell lines and for the HEK293T-trained Optimus 5-Prime
model against held-out measurements.}
\label{tab:cross_cell}
\begin{tabular}{lc}
\toprule
Comparison & $r^{2}$ \\
\midrule
HEK293T replicates & 0.938 \\
T cell replicates & 0.814 \\
HEK293T vs T cells (range) & 0.837--0.870 \\
HEK293T vs HepG2 (range) & 0.847--0.871 \\
\midrule
Optimus 5-Prime $\to$ HEK293T (held-out 20{,}000) & 0.937 \\
Optimus 5-Prime $\to$ T cells & 0.841 \\
Optimus 5-Prime $\to$ HepG2 & 0.840 \\
\bottomrule
\end{tabular}
\end{table}

\subsection{Designing high-MRL 50~nt 5$^\prime$UTRs with Fast SeqProp and DENs}

Having established that Optimus 5-Prime captures real biology, we next
generated synthetic 50~nt 5$^\prime$UTRs that maximize its predicted MRL.
We compared two generative strategies: gradient-based Fast
SeqProp~\citep{linder2021fastseqprop}, with and without VAE
regularization, and Deep Exploration Networks (DENs)~\citep{linder2020den}
(Figure~\ref{fig:den_schematic}). To validate that these generators
truly produced high-MRL sequences and were not simply gaming a single
predictor, we also trained an independent linear k-mer model as a
``second opinion'': it achieved Pearson $r = 0.5213$ on a held-out test
set, and it scored the Fast SeqProp-designed sequences within the top
1\% of its predictions on the entire test set. This convergent
validation from a structurally different model gave us confidence that
the designed sequences had real high-MRL features rather than
predictor-specific artifacts.

We then placed designed and control 50~nt 5$^\prime$UTRs upstream of
megaTAL mRNAs targeting \emph{TGFBR2} or \emph{PDCD1}, and
electroporated the resulting mRNAs at 0.25--2~pmol into K562 cells. As
expected, editing efficiency scaled with mRNA dose across all
5$^\prime$UTRs. At 2~pmol mRNA, several designs exceeded 40\% absolute
editing for the \emph{TGFBR2} megaTAL and 80\% absolute editing for the
\emph{PDCD1} megaTAL, levels that are well above the range typically
reported for low-dose mRNA editing. However, the gradient-based
strategy showed a telling failure mode: roughly 50\% of the Fast
SeqProp-designed sequences, regardless of VAE regularization,
delivered lower editing efficiencies than their high predicted MRL
would suggest. This discrepancy motivated a more rigorous comparison
on a different library length, which we describe next.

\subsection{A retrained model is required for 25~nt fully-variable 5$^\prime$UTRs}

The fixed-end 50~nt library holds the start context constant and allows
only the distal portion of the 5$^\prime$UTR to vary. For applications
that demand shorter, fully random 5$^\prime$UTRs, we sequenced a
random-end 25~nt library. Inter-replicate MRL correlation on this
library was lower than on the fixed-end 50~nt data:
$r^{2} = 0.692$ on the top 20{,}000 sequences by coverage and
$r^{2} = 0.844$ on the top 2{,}000, compared with $r^{2} = 0.938$ for
the 50~nt replicates. We therefore took the top-2{,}000 $r^{2}$ of
$0.844$ as an upper bound on how well any model could predict MRL on
this library.

To test whether previously developed Optimus variants could be reused,
we evaluated two HEK293T-trained models on the top 2{,}000 sequences
of the 25~nt library. The model trained on the fixed-end 50~nt library
achieved only $r^{2} = 0.564$, and an Optimus 5-Prime variant trained
on a 25--100~nt fixed-start library improved only marginally to
$r^{2} = 0.600$ (Figure~\ref{fig:model_25nt},
Table~\ref{tab:model_perf}). In other words, models trained on
fixed-start libraries transfer poorly to fully-variable short
5$^\prime$UTRs, consistent with the absence of a learned, length- and
context-matched start-codon grammar.

We therefore retrained the same architecture directly on the random-end
25~nt MPRA data, yielding a model we call Optimus 5-Prime(25). On the
same held-out test set, Optimus 5-Prime(25) reached $r^{2} = 0.806$,
within hailing distance of the inter-replicate ceiling of $0.844$. A
companion k-mer validator for the 25~nt setting reached Pearson
$r = 0.4094$ on held-out data and, as with the 50~nt setting, placed
designed sequences mostly within the top 1\% of its predictions on the
full test set, with a few outliers within the top 1.2\%, 6.3\%, or 4.5\%
--- all of which gave us usable confidence intervals on Fast SeqProp
and DEN designs at the 25~nt scale. The lesson is pragmatic: for
5$^\prime$UTRs shorter than the training library's start-context window,
retraining on length-matched data is necessary.

\begin{table}[htbp]
\centering
\caption{Performance of 5$^\prime$UTR models on the random-end 25~nt
MPRA test set.}
\label{tab:model_perf}
\begin{tabular}{lc}
\toprule
Model & $r^{2}$ on top-2{,}000 sequences \\
\midrule
Optimus 5-Prime (fixed-end 50~nt, ported) & 0.564 \\
Optimus 5-Prime 25--100~nt (fixed-start, ported) & 0.600 \\
Optimus 5-Prime(25) (this work) & 0.806 \\
\midrule
Inter-replicate ceiling (top 2{,}000) & 0.844 \\
\bottomrule
\end{tabular}
\end{table}

\subsection{Designed 5$^\prime$UTRs drive strong mRNA-delivered megaTAL editing}

We next asked whether designed 5$^\prime$UTRs translate into higher
gene editing efficiency in K562 cells, using megaTAL nucleases
targeting \emph{TGFBR2} and \emph{PDCD1} delivered as mRNA at 2~pmol.

For \emph{TGFBR2}, the best DEN-designed 5$^\prime$UTR produced
55.6\% absolute editing efficiency, an 18--33\% improvement over the
previous best natural control LAMA5 and the highest efficiency of any
5$^\prime$UTR tested at every mRNA dose (Figure~\ref{fig:editing},
Table~\ref{tab:editing}). Most of our designs matched or exceeded the
performance of 5$^\prime$UTRs drawn from the top 0.02\% of the random
MPRA libraries, and one design gave up to a 50\% higher absolute
\emph{TGFBR2} editing efficiency than any natural control. These
results indicate that generative design can indeed raise the ceiling
of mRNA-delivered editing activity in a practical therapeutic context.

The \emph{PDCD1} megaTAL, however, revealed an important wrinkle. The
same DEN-designed 5$^\prime$UTR that was best for \emph{TGFBR2}
delivered only 80.8\% absolute editing on \emph{PDCD1}, indistinguishable
from the Strong Kozak control, while the VAT1 natural control performed
best in this context with 91.4\% absolute editing
(Figure~\ref{fig:editing}, Table~\ref{tab:editing}). In other words,
the ranking of 5$^\prime$UTRs changed with cargo: the top MRL-predicted
design was not the top editor for every target. This cargo dependence,
together with the earlier observation that 50\% of
Fast SeqProp-designed sequences under-delivered relative to their
predicted MRL, indicates that predicted MRL is a useful but imperfect
proxy for functional output.

Taken as a whole, these data show that DEN-designed
5$^\prime$UTRs can exceed the best natural controls for specific
cargos, but that the 5$^\prime$UTR-cargo interaction is not captured
by a single MRL predictor and should still be validated empirically
in the deployment context.

\begin{table}[htbp]
\centering
\caption{Absolute megaTAL editing efficiency at 2~pmol mRNA in K562
cells for selected 5$^\prime$UTRs. The DEN-designed sequence is the
highest-MRL DEN 5$^\prime$UTR from this work; LAMA5 and VAT1 are
natural high-translation controls; ``Strong Kozak'' is a synthetic
canonical start-context control.}
\label{tab:editing}
\begin{tabular}{lcc}
\toprule
5$^\prime$UTR & \emph{TGFBR2} editing (\%) & \emph{PDCD1} editing (\%) \\
\midrule
DEN-designed (best on \emph{TGFBR2}) & 55.6 & 80.8 \\
LAMA5 (previous best, \emph{TGFBR2}) & 41.8--47.1$^{\dagger}$ & --- \\
VAT1 (best control on \emph{PDCD1}) & --- & 91.4 \\
Strong Kozak control & --- & 80.8$^{\ddagger}$ \\
\bottomrule
\end{tabular}
\\[2pt]
{\footnotesize
$^{\dagger}$LAMA5 range derived from the stated 18--33\% improvement
of the DEN-designed sequence over LAMA5 at 2~pmol.
$^{\ddagger}$Text states the DEN-designed sequence performed
``only as well as the Strong Kozak control'' on \emph{PDCD1}.
Dashes indicate cargo/control combinations not explicitly numeric in
the source data.}
\end{table}

\section{Discussion}

The experiments presented above make three points about model-guided
5$^\prime$UTR design for mRNA therapeutics. First, 5$^\prime$UTR
activity is portable across human cell types. Second, existing
5$^\prime$UTR models do not transfer across library geometries and
must be retrained for fully-variable short UTRs. Third, generatively
designed 5$^\prime$UTRs can outperform the best natural controls in
mRNA-delivered gene editing, but this superiority is cargo- and
locus-specific. Each of these points has practical consequences for
how one should integrate deep learning into an mRNA therapeutic
discovery pipeline.

The observation that inter-cell-line MRL correlation
($r^{2} = 0.837$--$0.871$) approaches the inter-replicate ceiling
within a single cell line ($r^{2} = 0.814$--$0.938$) places a useful
upper bound on how much cell-type-specific 5$^\prime$UTR engineering
can actually help in the cell lines surveyed here. The bulk of
5$^\prime$UTR activity reflects conserved translational machinery
shared across HEK293T, T cells, and HepG2, and the HEK293T-trained
Optimus 5-Prime model's near-linear drop from $r^{2} = 0.937$ in
HEK293T to $r^{2} = 0.841$ and $0.840$ in T cells and HepG2 follows
exactly the inter-cell-line noise floor. This is consistent with work
on vertebrate 5$^\prime$UTR grammar~\citep{reimaopinto2025danio5p,
strayer2024naptrap} and on Kozak-context-driven translation
control~\citep{chaldebas2026utrvariants,stettler2026kozak}. In a
practical design loop, it means a single high-quality training set is
often sufficient to seed design in several target cell types, provided
one accepts a modest and predictable loss in predictive accuracy.

The failure of Optimus 5-Prime models trained on fixed-end 50~nt or
25--100~nt libraries to generalize to a fully-variable 25~nt library
($r^{2} = 0.564$ and $r^{2} = 0.600$, versus $0.806$ for the retrained
Optimus 5-Prime(25)) is consistent with the increasingly clear picture
that short 5$^\prime$UTRs place the start codon in a distinct scanning
regime~\citep{chaldebas2026utrvariants,balasov2025kozaknoncanonical,
shikalov2021ddx3}. In a 25~nt fully-variable UTR, there is no fixed
``scanning tract,'' and the proximity of uAUGs and non-canonical
initiation codons to the cap becomes a first-order determinant of
output. This favors library-matched models and argues against the
common practice of reusing a single, often HEK293T-based, 5$^\prime$UTR
model across unrelated UTR geometries in the mRNA therapeutics literature.

Perhaps the most therapeutically consequential finding is that the
DEN-designed 5$^\prime$UTR that was best for the \emph{TGFBR2} megaTAL
(55.6\% absolute editing at 2~pmol mRNA; 18--33\% improvement over
LAMA5) was \emph{not} the best 5$^\prime$UTR for the \emph{PDCD1}
megaTAL (80.8\% editing, matching the Strong Kozak control), where the
VAT1 natural UTR reached 91.4\%. Several, non-exclusive, explanations
are plausible: the cargo coding sequence can interact with the
5$^\prime$UTR to modulate secondary structure around the start codon;
different megaTAL sequences may have different sensitivity to absolute
enzyme concentration; and the underlying translation machinery balance
may differ subtly at the two genomic loci. Whatever the mechanism,
the outcome is that 50\% of Fast SeqProp-designed sequences showed
lower editing efficiency than their predicted MRL would suggest, and
the top MRL predictor does not necessarily select the top editor for
every cargo. Integrating functional readouts --- here, absolute
editing efficiency --- directly into the design loop is therefore
essential for translating 5$^\prime$UTR modelling into therapeutic
gains, echoing recent work that couples UTR tuning to formulation
and delivery~\citep{kim2026lnp_utr,phua2026ns1translation}.

Several limitations of the present study should be noted honestly. We
examined two gene targets (\emph{TGFBR2}, \emph{PDCD1}) and two cell
lines (K562 and, by reference, HepG2) under a limited mRNA dose range
(0.25--2~pmol), which is appropriate for proof-of-principle but not
for predicting clinical behavior. The nuclease modality studied here
is megaTALs~\citep{gwiazda2016tcellediting,mcmurrough2018laglidadg,
laforet2019endonuclease}, and whether our 5$^\prime$UTRs transfer to
CRISPR/Cas9 or base-editor payloads has not been tested, although the
principle of MRL-guided design should be payload-agnostic. Our MPRAs
were performed in immortalized lines and in primary T cells, not in
the target cells most relevant for clinical adoptive cell
therapies~\citep{fix2022tgfbr2,mohammad2025pd1}, and the proportion of
MRL variance that reflects cell-type-specific RNA-binding-protein
biology~\citep{shikalov2021ddx3} rather than core translational
machinery is likely to grow as one moves from HEK293T and K562
towards primary effector T cells \emph{in vivo}.

Looking forward, three directions are immediate. First, coupling
library-matched 5$^\prime$UTR models to the formulation and
chemistry advances documented in recent LNP and m$^{1}\Psi$
work~\citep{tang2023smart5utr,kim2026lnp_utr,thomas2021bnt162b2}
should compound gains across the mRNA stack. Second, training
objectives for generative models like
DENs~\citep{linder2020den,linder2021fastseqprop} can be extended to
optimize cargo-specific editing efficiency directly, rather than
indirectly through predicted MRL, closing the loop between
\emph{in silico} fitness and functional output. Third, because the
5$^\prime$UTR is a natural locus for cargo-specific fine tuning ---
cheap to vary, orthogonal to the payload sequence, and directly tied
to dose --- it is a good candidate for personalized mRNA design
workflows of the kind now being piloted in
individualized cancer vaccines~\citep{sahin2026tnbcvaccine}. More
broadly, our results argue that model-guided 5$^\prime$UTR design
is a practical tuning knob for mRNA therapeutics, but one that
only pays off when coupled to empirical validation in the
deployment context.

\section*{Data and Code Availability}

All experimental data described here are summarized in the statistical
reports and methods sections reproduced verbatim from the authors'
laboratory notes. Model code is available at the public repository
associated with the experimental log.

\bibliography{references}

\end{document}
